\slajdid{04-s006}
\begin{frame}[label=04-s006]{Primer: iz osnove 7 u osnovu 10}
% element: 04-s006:e01 — Primer 326.23 u osnovi 7 i sva težinska množenja; 167+17/49
% element: 04-s006:e02 — Oba decimalna niza iz originala za 17/49 i kompletnu vrednost
% element: 04-s006:e03 — Zaključak o konačnom ili beskonačnom zapisu pri promeni osnove
\setlength{\abovedisplayskip}{4pt}\setlength{\belowdisplayskip}{4pt}
\[326.23_7={?}_{10}\]
\[\begin{aligned}326.23_7&=3\cdot7^2+2\cdot7^1+6\cdot7^0+2\cdot7^{-1}+3\cdot7^{-2}\\&=167+\frac{17}{49}\end{aligned}\]
Želimo standardan pozicioni zapis, umesto zapisa razlomkom.
\[\frac{17}{49}=0.346938775510204081632653061224489795\ldots\]
\[326.23_7=167.346938775510204081632653061224489795\ldots_{10}\]
\Rezultat{Konačan broj cifara u jednom sistemu ne mora dati konačan broj cifara u drugom. Možda hoće, možda neće.}
\end{frame}
